% Compile with XeLaTeX. Content converted from semestrova-kontrolna-stem-5-klas.docx.
\documentclass[11pt,a4paper]{article}
\usepackage{fontspec}
\usepackage[ukrainian]{babel}
\usepackage[margin=18mm]{geometry}
\usepackage{array,tabularx}
\setmainfont{Arial}
\setlength{\parindent}{0pt}
\setlength{\parskip}{6pt}
\renewcommand{\arraystretch}{1.65}
\begin{document}
\begin{center}\Large\textbf{Семестрова контрольна робота з курсу STEM}\end{center}

5 клас\par

Укладено на основі навчального плану модулів «Я у Всесвіті», «Я так бачу!» і «Під знаком STEM».\par

Прізвище, ім'я \rule{0.65\linewidth}{0.4pt}\par

Клас \rule{0.25\linewidth}{0.4pt}    Дата \rule{0.25\linewidth}{0.4pt}\par

\section*{Варіант 1}

\subsection*{1. Тестові завдання}

1. Який прилад допомагає спостерігати небесні тіла?\newline А. Термометр\newline Б. Телескоп\newline В. Мікрофон\newline Г. Компас\par

2. Що належить до сучасних візуальних медіа?\newline А. Анімація\newline Б. Формула\newline В. Таблиця множення\newline Г. Лінійка\par

3. Який із наведених прикладів є штучним знаком?\newline А. Темні хмари\newline Б. Слід лапи на снігу\newline В. Дорожній знак\newline Г. Веселка\par

\subsection*{2. Теоретичне питання}

Поясни, яку роль у житті людини відіграють зображення, символи та знаки. Напиши 5-7 речень.\par

\rule{0.96\linewidth}{0.4pt}\par

\rule{0.96\linewidth}{0.4pt}\par

\rule{0.96\linewidth}{0.4pt}\par

\rule{0.96\linewidth}{0.4pt}\par

\rule{0.96\linewidth}{0.4pt}\par

\newpage\textbf{Варіант 1 — продовження}\par

\subsection*{3. Практичне завдання 1}

Розподіли подані поняття у три групи: «Космос», «Графіка», «Символи». Запиши всі слова в таблицю.\par

Слова для розподілу: телескоп, герб, фотографія, орбіта, емблема, анімація, ракета, піктограма, планета, малюнок, печатка, колаж.\par

\begin{tabularx}{\linewidth}{|>{\raggedright\arraybackslash}X|>{\raggedright\arraybackslash}X|>{\raggedright\arraybackslash}X|}\hline

\textbf{Космос} & \textbf{Графіка} & \textbf{Символи} \\ \hline

 &  &  \\ \hline

 &  &  \\ \hline

 &  &  \\ \hline

 &  &  \\ \hline

\end{tabularx}\par

\subsection*{4. Практичне завдання 2}

Заповни таблицю «Об'єкт - призначення - модуль».\par

\begin{tabularx}{\linewidth}{|>{\raggedright\arraybackslash}X|>{\raggedright\arraybackslash}X|>{\raggedright\arraybackslash}X|}\hline

\textbf{Об'єкт} & \textbf{Для чого використовується} & \textbf{До якого модуля належить} \\ \hline

Телескоп &  &  \\ \hline

Герб &  &  \\ \hline

Анімація &  &  \\ \hline

Ракета &  &  \\ \hline

Піктограма &  &  \\ \hline

Фотографія &  &  \\ \hline

Печатка &  &  \\ \hline

Орбіта &  &  \\ \hline

\end{tabularx}\par

\newpage

\section*{Варіант 2}

\subsection*{1. Тестові завдання}

1. Що є небезпечним фактором у відкритому космосі?\newline А. Космічне сміття\newline Б. Дощ\newline В. Туман\newline Г. Пісок\par

2. Що з наведеного найбільше пов'язане з комп'ютерною графікою?\newline А. Створення зображень на комп'ютері\newline Б. Вирощування рослин\newline В. Вимірювання температури\newline Г. Приготування їжі\par

3. Який приклад є природним знаком?\newline А. Смайлик\newline Б. Літера\newline В. Темні хмари\newline Г. Печатка\par

\subsection*{2. Теоретичне питання}

Поясни письмово, що таке знакова система та як знаки допомагають людям передавати інформацію. Напиши 5-7 речень.\par

\rule{0.96\linewidth}{0.4pt}\par

\rule{0.96\linewidth}{0.4pt}\par

\rule{0.96\linewidth}{0.4pt}\par

\rule{0.96\linewidth}{0.4pt}\par

\rule{0.96\linewidth}{0.4pt}\par

\newpage\textbf{Варіант 2 — продовження}\par

\subsection*{3. Практичне завдання 1}

Заповни таблицю «Ризик і захист у космосі».\par

\begin{tabularx}{\linewidth}{|>{\raggedright\arraybackslash}X|>{\raggedright\arraybackslash}X|>{\raggedright\arraybackslash}X|}\hline

\textbf{Небезпечний фактор} & \textbf{Чим загрожує} & \textbf{Спосіб захисту} \\ \hline

Холод &  &  \\ \hline

Космічне сміття &  &  \\ \hline

Радіація &  &  \\ \hline

\end{tabularx}\par

\subsection*{4. Практичне завдання 2}

Склади план створення інформаційного плаката «Космічні професії майбутнього». У плані має бути не менше 5 кроків.\par

1. \rule{0.65\linewidth}{0.4pt}\par

2. \rule{0.65\linewidth}{0.4pt}\par

3. \rule{0.65\linewidth}{0.4pt}\par

4. \rule{0.65\linewidth}{0.4pt}\par

5. \rule{0.65\linewidth}{0.4pt}\par

6. \rule{0.65\linewidth}{0.4pt}\par

\newpage

\section*{Тестовий варіант-зразок для учнів}

Цей варіант можна використати як тренувальний перед контрольною роботою.\par

\subsection*{1. Тестові завдання}

1. Який прилад допомагає спостерігати небесні тіла?\newline А. Термометр\newline Б. Телескоп\newline В. Мікрофон\newline Г. Компас\par

2. Що належить до сучасних візуальних медіа?\newline А. Анімація\newline Б. Таблиця множення\newline В. Лінійка\newline Г. Формула\par

3. Який із прикладів є штучним знаком?\newline А. Веселка\newline Б. Темні хмари\newline В. Дорожній знак\newline Г. Слід лапи на снігу\par

\subsection*{2. Теоретичне питання}

Поясни, як люди використовують знаки й зображення для передавання інформації.\par

\rule{0.96\linewidth}{0.4pt}\par

\rule{0.96\linewidth}{0.4pt}\par

\rule{0.96\linewidth}{0.4pt}\par

\rule{0.96\linewidth}{0.4pt}\par

\rule{0.96\linewidth}{0.4pt}\par

\newpage\textbf{Тестовий варіант-зразок для учнів — продовження}\par

\subsection*{3. Практичне завдання 1}

Розподіли слова у три колонки: «Космос», «Графіка», «Символи».\par

Слова: ракета, малюнок, печатка, планета, фотографія, знак, телескоп, емблема, орбіта, анімація, герб, піктограма.\par

\begin{tabularx}{\linewidth}{|>{\raggedright\arraybackslash}X|>{\raggedright\arraybackslash}X|>{\raggedright\arraybackslash}X|}\hline

\textbf{Космос} & \textbf{Графіка} & \textbf{Символи} \\ \hline

 &  &  \\ \hline

 &  &  \\ \hline

 &  &  \\ \hline

 &  &  \\ \hline

\end{tabularx}\par

\subsection*{4. Практичне завдання 2}

Заповни таблицю «Приклад - група - коротке пояснення».\par

\begin{tabularx}{\linewidth}{|>{\raggedright\arraybackslash}X|>{\raggedright\arraybackslash}X|>{\raggedright\arraybackslash}X|}\hline

\textbf{Приклад} & \textbf{До якої групи належить} & \textbf{Коротке пояснення} \\ \hline

Телескоп &  &  \\ \hline

Герб &  &  \\ \hline

Анімація &  &  \\ \hline

Планета &  &  \\ \hline

Піктограма &  &  \\ \hline

Фотографія &  &  \\ \hline

\end{tabularx}\par
\end{document}
